% !TEX program = xelatex
\documentclass[11pt]{article}
\usepackage{fontspec}
\usepackage{xeCJK}
\setmainfont{CMU Serif}
\IfFontExistsTF{Noto Serif CJK SC}{\setCJKmainfont{Noto Serif CJK SC}}{%
  \IfFontExistsTF{Source Han Serif SC}{\setCJKmainfont{Source Han Serif SC}}{%
    \setCJKmainfont{FandolSong-Regular}%
  }%
}%
\usepackage[margin=1in]{geometry}
\usepackage{amsmath}
\usepackage{amssymb}
\usepackage{booktabs}
\usepackage{siunitx}
\usepackage{array}

\sisetup{table-number-alignment=center}

\setlength{\parindent}{0pt}
\setlength{\parskip}{0.35\baselineskip}
\setlength{\tabcolsep}{4pt}
\renewcommand{\arraystretch}{0.92}

\makeatletter
\renewcommand{\maketitle}{%
  \begin{center}
    {\LARGE \@title\par}
  \end{center}
  \vspace{-0.6\baselineskip}
}
\makeatother

\title{Idea报告草稿：MRC-SID}
\author{}
\date{}

\begin{document}

\maketitle

\section{研究背景}

生成式推荐将目标 item 的预测改写为离散 token 序列生成问题。以 MiniOneRec 为代表的方法通常先为每个 item 构造固定长度的语义 ID（Semantic ID, SID），再用大语言模型根据用户历史自回归生成目标 item 的 SID。该路线的优点是能够复用语言模型的训练与解码范式，但也引入了一个新的核心问题：\textbf{tokenizer 的设计会直接影响后续生成式推荐的可学习性与可预测性。}

在当前主流 SID 路线中，一个非常常见的设计是：先从 item 的文本语义或语义增强表征出发，再将其整体量化为多层离散 code。最近已有不少工作开始将 collaborative signal 注入 tokenizer，但多数方法仍然把协同信息视为一个\textbf{统一资源}，并尝试在整个 tokenizer 中做全局融合。这样的设计虽然自然，却未必与 hierarchical SID 的结构相匹配。

在我们的仓库中，当前最强基线已经能够较稳定地复现 MiniOneRec 主链，并在 Industrial 与 Office 上接近论文结果。因此，当前真正值得研究的问题已经不再是“系统能否跑通”，而是：\textbf{在当前 SID 设计下，性能瓶颈究竟在哪里，以及 collaborative information 应该以什么粒度进入 hierarchical tokenizer。}

\section{当前观察与问题诊断}

\subsection{全局 collision 存在，但不是主矛盾}

基于当前最优 Industrial 与 Office 结果的诊断，我们发现：完整 SID collision 的确存在，但绝对比例并不高。Industrial 与 Office 的 collision rate 都约为 $0.43\%$。这说明“不同 item 完全撞到同一个 SID”不是不存在，但不足以解释当前主要性能瓶颈。

因此，当前问题不宜被表述为“global collision 太严重”，而更应被理解为：\textbf{SID 结构虽然整体可用，但在局部层级上的判别性仍然不足。}

\subsection{主要错误模式是 local leaf ambiguity}

进一步分析 evaluate 结果后，我们观察到：
\begin{enumerate}
  \item 模型并不是经常跑到完全错误的语义子树；
  \item 相反，很多错误样本已经写对了第一层，甚至前两层 prefix；
  \item 最终错误主要集中在最后一层 leaf token 的选择上。
\end{enumerate}

这意味着当前更像是如下问题：
\begin{quote}
模型已经走到了大致正确的 prefix 子空间，但在同一个局部子树内部，仍然难以可靠地区分多个语义邻近、文本相似、但行为上应当区分的 item。
\end{quote}

换言之，当前主要瓶颈不是 global routing，而是 \textbf{local leaf ambiguity}。

\subsection{协同信号确实缺失，但其作用不是单尺度的}

我们进一步做了 train-only 的协同兼容性 probe：给定测试样本历史，分别用不同长度的历史窗口构造协同得分，比较真实 target 与预测 SID 下最强候选之间的协同兼容性。这里考虑四种 view：
\begin{itemize}
  \item \texttt{coarse10}: 最近 10 个历史 item；
  \item \texttt{mid3}: 最近 3 个历史 item；
  \item \texttt{local2}: 最近 2 个历史 item；
  \item \texttt{fine1}: 仅最近 1 个历史 item。
\end{itemize}

Pilot 结果表明，不同 collaborative view 在不同错误桶上的作用并不相同。Industrial 与 Office 上都出现了类似趋势：
\begin{itemize}
  \item \texttt{coarse10} 在整体错误池上最强或接近最强；
  \item \texttt{mid3} 在 deepest local ambiguity（same-$l_2$）上最强；
  \item \texttt{fine1} 最局部，但也最稀疏、最脆弱。
\end{itemize}

这说明当前不能再把 collaborative information 理解成一个统一向量，然后在整个 hierarchical SID 上均匀使用。更合理的假设是：
\begin{quote}
不同 SID level 可能需要不同分辨率的 collaborative signal，而且这种分配关系应当被学习，而不是被手工指定。
\end{quote}

\begin{table}[t]
  \centering
  \caption{不同 collaborative view 在不同错误桶上的 target-better rate（节选）。}
  \small
  \begin{tabular}{l l S[table-format=1.4] S[table-format=1.4]}
    \toprule
    数据集 & View & {All Errors} & {Same-$l_2$ Errors} \\
    \midrule
    Industrial & \texttt{coarse10} & 0.0981 & 0.2284 \\
    Industrial & \texttt{mid3}     & 0.0888 & 0.2346 \\
    Industrial & \texttt{fine1}    & 0.0776 & 0.2006 \\
    \midrule
    Office     & \texttt{coarse10} & 0.1083 & 0.4048 \\
    Office     & \texttt{mid3}     & 0.1072 & 0.4167 \\
    Office     & \texttt{fine1}    & 0.0836 & 0.3690 \\
    \bottomrule
  \end{tabular}
\end{table}

\section{Motivation}

基于以上观察，我们认为当前最值得提出的问题不是：
\begin{itemize}
  \item “是否应将协同信息注入 tokenizer”，
\end{itemize}
而是：
\begin{quote}
\textbf{在 hierarchical SID tokenizer 中，不同 level 应该如何分配不同分辨率的 collaborative views？}
\end{quote}

这个问题的关键在于，hierarchical SID 的不同层级本身承担着不同角色：
\begin{itemize}
  \item 上层 token 更接近粗粒度 routing；
  \item 中层 token 更接近 subgroup organization；
  \item 深层 token 更接近 leaf-level discrimination。
\end{itemize}

与此同时，不同 collaborative signal 也携带不同的稳定性与噪声特性：
\begin{itemize}
  \item 更 coarse 的 view 覆盖更广、统计更稳定；
  \item 更 local 的 view 更贴近瞬时转移，但也更 sparse、更 noisy。
\end{itemize}

因此，一个合理的 hierarchical SID tokenizer 不应在所有层级上使用同一 collaborative representation。相反，它应当学习：
\begin{itemize}
  \item 哪些 level 更适合偏稳定的 coarse signal；
  \item 哪些 level 更适合偏局部的 mid/local signal；
  \item 这种分配是否真的非均匀。
\end{itemize}

\section{方法：MRC-SID}

\subsection{核心思想}

我们提出 \textbf{MRC-SID（Multi-Resolution Collaborative Allocation for Hierarchical Semantic IDs）}。

该方法的核心思想是：
\begin{enumerate}
  \item 从 train-only interaction data 构造多种分辨率的 collaborative views；
  \item 对不同 collaborative view 分别做净化与去噪；
  \item 在 hierarchical SID tokenizer 的每一层，学习一个独立的 collaborative allocation；
  \item 使不同 SID level 自动学习自己更适合的 semantic/collaborative mixture。
\end{enumerate}

因此，MRC-SID 不是“再做一个 uniform collaborative fusion”，而是在问：
\begin{quote}
当 tokenizer 本身是 hierarchical 的时候，collaborative information 是否也应该被 hierarchy-aware 地组织与分配？
\end{quote}

\subsection{Multi-Resolution Collaborative View Bank}

对于每个 item $i$，我们首先从 train-only 交互数据中构造三种 collaborative view：
\begin{itemize}
  \item \textbf{coarse view} $\mathbf{c}_i^{(c)}$：长窗口、稳定、覆盖广；
  \item \textbf{mid view} $\mathbf{c}_i^{(m)}$：中尺度上下文结构；
  \item \textbf{local view} $\mathbf{c}_i^{(l)}$：短程 recent transition signal。
\end{itemize}

与 naive 的 single-view 方案不同，这三种表示不是简单重复，而是刻意对应不同的分辨率与噪声 profile。

\subsection{View-Specific Purification}

考虑到我们已经在当前仓库中观察到 naive global fusion 会导致结构性塌缩，因此对于每个 collaborative view，我们都引入独立的 purification 流程：
\begin{itemize}
  \item coarse view 采用更偏平滑、降噪、低秩压缩的处理；
  \item mid view 采用更偏局部结构过滤的处理；
  \item local view 采用更激进的 support pruning 与 recency-aware smoothing。
\end{itemize}

记净化后的三种 collaborative representation 为：
\[
\tilde{\mathbf{c}}_i^{(c)},\quad
\tilde{\mathbf{c}}_i^{(m)},\quad
\tilde{\mathbf{c}}_i^{(l)}.
\]

与此同时，记 item 的 semantic representation 为 $\mathbf{s}_i$。

\subsection{Level-Wise Collaborative Allocation}

对于 3-level SID，我们为每一层 $\ell \in \{1,2,3\}$ 学习一个独立 gate：
\[
\boldsymbol{\alpha}^{(\ell)} =
\mathrm{softmax}\!\left(\mathbf{w}^{(\ell)}\right)
\in \mathbb{R}^{4},
\]
其中四个分量分别对应：
\[
\mathbf{s}_i,\quad
\tilde{\mathbf{c}}_i^{(c)},\quad
\tilde{\mathbf{c}}_i^{(m)},\quad
\tilde{\mathbf{c}}_i^{(l)}.
\]

于是，第 $\ell$ 层的融合表示可以写为：
\[
\mathbf{z}_i^{(\ell)} =
\alpha_{\text{sem}}^{(\ell)} \mathbf{s}_i
 \alpha_{c}^{(\ell)} \tilde{\mathbf{c}}_i^{(c)}
 \alpha_{m}^{(\ell)} \tilde{\mathbf{c}}_i^{(m)}
 \alpha_{l}^{(\ell)} \tilde{\mathbf{c}}_i^{(l)}.
\]

这一设计有两个重要含义：
\begin{enumerate}
  \item 我们不再手工写死“上层一定 coarse、下层一定 local”；
  \item 我们允许模型自己发现不同 level 的最优 collaborative mixture。
\end{enumerate}

\subsection{Hierarchy-Aware Tokenization}

在得到各层专属的融合表示 $\mathbf{z}_i^{(1)}, \mathbf{z}_i^{(2)}, \mathbf{z}_i^{(3)}$ 后，再将它们送入当前 hierarchical SID tokenizer 对应层的量化与分配模块，从而得到最终的 3-level SID：
\[
\text{SID}(i) = \left(q_i^{(1)}, q_i^{(2)}, q_i^{(3)}\right).
\]

这里的关键并不是重新定义整个 tokenizer 框架，而是把 collaborative allocation 明确建模为一个 \textbf{SID-level principle}。

\section{为什么这条方法与已有方向不同}

MRC-SID 与已有思路的差别主要体现在三个方面：

\subsection{不同于 generic front-end collaborative tokenizer}

已有很多方法讨论“如何把 collaboration 加进 tokenizer”，但它们往往把 collaborative signal 当成整体融合资源。MRC-SID 想证明的是：
\begin{quote}
collaborative information 的 utility 本身是 hierarchy-dependent 的。
\end{quote}

\subsection{不同于 post-hoc local repair}

像 ACLR-lite 这样的后处理方法能够证明 local repair 有 headroom，但其核心仍然停留在 inference-time refinement。MRC-SID 则试图把 collaborative structure 更早地推进到 tokenizer 内部，而且不是 uniform 地推进，而是按 level 分配。

\subsection{不同于固定 heuristic mapping}

MRC-SID 的主张不是：
\begin{quote}
Level 1 必须 coarse，Level 3 必须 local。
\end{quote}

MRC-SID 的主张是：
\begin{quote}
每一层的 collaborative allocation 应当被学习出来，因为 uniform fusion 与 fixed heuristic mapping 都可能是结构错配。
\end{quote}

\section{可证伪主张}

这条方法线的核心不是“肯定有效”，而是非常明确、可证伪：
\begin{enumerate}
  \item MRC-SID 应优于 semantic-only tokenizer；
  \item MRC-SID 应优于 parameter-matched 的 uniform all-view fusion；
  \item MRC-SID 应优于固定的 coarse/mid/local hand-designed assignment；
  \item learned gate 在各层应当呈现可解释的 non-uniform allocation；
  \item deeper levels 应更可能受益于 non-global views。
\end{enumerate}

如果这些现象不成立，那么这条方法应被降级为一个分析性 insight，而不是完整 paper method。

\section{最小实验路线}

为了避免再次把方法做大做杂，MRC-SID 的第一版实现应该保持克制：
\begin{itemize}
  \item 只保留三种 collaborative view；
  \item 每个 view 只保留一套最小 purification；
  \item 每一层只保留一个 learned gate；
  \item 暂不引入 ambiguity-aware dynamic gating；
  \item 先只在 Industrial 上做主实验。
\end{itemize}

最小对照应至少包括：
\begin{enumerate}
  \item semantic-only baseline；
  \item uniform all-view fusion；
  \item fixed level assignment；
  \item MRC-SID learned level allocation。
\end{enumerate}

真正决定这条路线是否值得继续的，不是“是否比 semantic-only 稍好”，而是：
\begin{quote}
\textbf{MRC-SID 是否能稳定优于 uniform fusion，并且学出非均匀、可解释的 level-wise allocation。}
\end{quote}

\section{总结}

当前仓库中的观察已经表明：
\begin{itemize}
  \item global collision 不是主矛盾；
  \item local leaf ambiguity 才是更核心的失败模式；
  \item collaborative signal 对不同错误深度的帮助并不相同；
  \item naive global fusion 会导致 tokenizer collapse。
\end{itemize}

因此，当前更合理的问题不再是“是否要加 collaborative information”，而是：
\begin{quote}
\textbf{如何在 hierarchical SID 中，将不同分辨率的 collaborative signal 按 level 合理分配。}
\end{quote}

MRC-SID 正是在这一问题上提出的一个更结构化、更可论文化的回答。它试图把 collaborative signal 从“统一融合资源”改写为“需要按 SID level 学习分配的层级资源”。如果后续实验能够证明 learned allocation 稳定优于 uniform fusion 与 fixed mapping，那么这条路线就有望成为当前这条前端 tokenizer 主线中最值得继续推进的方法。

\end{document}
